\section{Equal observed laws can imply different intervention answers}
\label{v13:la-equivalence}

A diagram assistant must recognize when available observations cannot determine
the requested direction or effect. Consider two project-constructed binary
models. Let $U$ be fair Bernoulli and $A=U$. In model one $Y=A$; in model two
$Y=U$. Both observed datasets contain only $(A,Y)=(0,0)$ and $(1,1)$ with equal
probability. Under an intervention setting $A=a$, model one has mean outcome
$a$ while model two has mean outcome $1/2$. Their average effects are one and
zero. Exact enumeration establishes the discrepancy.\cite{v13_project_exact}

Estimating the observed marginal propensity as $1/2$ does not solve the
problem: exchangeability given the available information is precisely what is
missing. A model that recognizes a plausible causal story may prefer one graph,
but that preference is not identification from these data. Expert knowledge,
measurements or experiments could distinguish the worlds; those inputs need
their own evidence and should remain visible in the result.

A continuous version shows why effect sets can be unbounded. Let independent
standard Gaussian variables $U$ and $\varepsilon$ generate
\[
 X=U,\qquad Y=bX-bU+\varepsilon.
\]
The factual law is $X=U,Y=\varepsilon$ for every real $b$. Under
$\operatorname{do}(X=x)$, however, the mean of $Y$ is $bx$. Therefore every
real intervention slope is compatible with the same observed distribution
within this allowed model class. More observations estimate that common law
more accurately; they do not select $b$. A discovery tool or language model
cannot remove this ambiguity without an additional restriction.

This is a useful evaluation case for agentic discovery. A high-quality response
retains the competing models, explains their observational equivalence and
identifies the information that could separate them. An overconfident single
graph should be scored as a substantive error even if its prose is convincing.
A refusal or whole-line effect set can be the scientifically correct output;
its lack of immediate decision usefulness should be reported separately.

\section{Prior knowledge needs a typed evidence record}

The inspected DirectLiNGAM interface distinguishes a directed path, no directed
path and unspecified prior knowledge. The y0 identification dispatcher requires
a supplied graph rather than eliciting its truth. These interfaces make the
translation from narrative evidence to formal restrictions consequential.
An agent saying ``A is likely relevant to Y'' has not supplied a justified
directed-path constraint.\cite{v13_sw_lingam,v13_sw_y0}

For each proposed restriction, a useful record distinguishes a temporal order,
an intervention mechanism, a measured association, a domain assertion and a
formal exclusion. It binds the variables, their measurement/availability clocks,
the source passage, the supporting scope and any disagreement. A source about
one population or intervention need not justify a hard exclusion in another.
The record should preserve which conclusions change if the assertion is removed.

Suppose a manager states that a rate change cannot be caused by next month's
withdrawal. That supports a temporal ordering when the rate decision was fixed
before the withdrawal. It does not rule out an earlier expectation of withdrawal
affecting the rate, or a common market shock affecting both. A diagram assistant
can expose the omitted expectation variable and the common-cause alternative.
It should not convert the temporal statement into unconfoundedness.

An analogous issue arises in network diagrams. MossSpider's selected source
uses a directed influence convention that requires explicit checking, while a
bank may draw arrows from guarantor to beneficiary. A language model can produce
the correct-looking graph but an incorrectly oriented adjacency representation.
A two-node semantic test with known distinct attributes is more informative
than a visual similarity score.\cite{v13_sw_mossspider}

\section{Set-valued reasoning should survive verbal reporting}

If uncertainty concerns several graphs or nuisance parameters, the effect
object may be a union of intervals. A project ratio example illustrates the
boundary. Let a numerator lie in $[1,2]$ and an audit multiplier in
$[-1/2,1/2]$, with the relation $m=r\theta$. Excluding division by zero while
retaining all feasible nonzero multipliers gives
\[
 \Theta=(-\infty,-2]\ \cup\ [2,\infty).
\]
The gap around zero and the unbounded tails are substantive. Replacing this
object by its hull loses the exclusion of zero; replacing it by a finite
interval loses compatible effects. A monetary transformation should be applied
to each component and preserve endpoint inclusion and infinite endpoints.

The agent should describe why the set is disconnected, what assumptions support
it and what information would make it more useful. It should not translate
``wide uncertainty'' into a invented central estimate. When both zero numerator
and zero multiplier remain jointly feasible, the effect can instead be the
whole real line. The distinction depends on the joint region, not just
separate point estimates.

This reasoning also separates model plausibility from sampling coverage. If a
joint confidence region contains the true observable and nuisance parameters,
its effect image contains the true effect on that event. The implication is
elementary; constructing the valid joint region under the bank's actual
dependence and audit law is the difficult statistical work. An agent can
propagate a supplied region faithfully. It cannot infer that its coverage is
valid merely because the downstream set arithmetic passed a test.
